The direct literature on seasonal variation in imaging-defined cardiac structure and function is sparse. The clearest direct evidence comes from repeated echocardiographic studies in athlete cohorts. Fagard et al. reported seasonal changes in cyclists, Crouse et al. described seasonal remodeling in women basketball athletes, Dernovoj et al. showed seasonal changes in right-heart functional measures in elite skiers, and Forsythe et al. observed that conventional structure and function were relatively stable across a rugby league season while selected speckle-tracking mechanics varied.

Outside athlete imaging, the literature becomes less direct. Human studies more consistently show seasonal variation in cardiovascular hemodynamics, autonomic measures, blood pressure variability, arterial stiffness, and heart-failure burden than in cardiac morphology itself. Winter is repeatedly associated with higher vascular resistance, altered cardiac output or stroke volume, elevated blood pressure variability, increased arterial stiffness, and higher rates of heart-failure hospitalization and mortality.

Taken together, the literature supports biological plausibility for seasonality in cardiac phenotypes, but direct evidence in general-population cardiac MRI remains very limited. This gap is important: it means that a well-controlled cardiac MRI analysis has the potential to provide genuinely novel evidence rather than merely replicating an already mature literature.

Among the currently retrieved open-access full texts, the most directly relevant imaging paper is Forsythe et al., where seasonal/training-phase variation appeared mainly in advanced mechanics rather than conventional chamber measures. This suggests that any seasonal MRI signal may be subtle and more likely to emerge in sensitive phenotypes such as strain, wall mechanics, or shape descriptors than in gross chamber size alone.
